\section{A calculus that records the norm of each observation}\label{sec:foundation}

\subsection{Arithmetic, coding, and the meaning of a presentation}

We use intuitionistic first-order logic with equality, the natural numbers,
a symbol for every fixed primitive-recursive definition, its defining
equations, successor injectivity and nonzeroness, and induction for
quantifier-free formulas. This is the repository's ordinary
\texttt{PRAAx} presentation \cite{repo-pr}. No induction over arbitrary
quantified statements is added.

A signed integer is a difference of two naturals. A rational is an integer
numerator and a positive natural denominator; it can be normalized by the
PR gcd algorithm but normalization is not needed for order. Equality of
$a/b$ and $c/d$ means $ad=bc$, and order means $ad\le bc$, with $b,d>0$.
A finite list is coded by the repository's pairing and sequence operations.
Every access is explicitly guarded by its stored length. Finite vectors,
matrices, Boolean tables, and polynomials are lists with dimension and
well-formedness checks. None requires a new arithmetic sort.

For clarity, a map displayed as $a:\N\to\Q$ below denotes a \emph{fixed
program schema} and its numerical parameters. It does not denote an
arbitrary function variable in the language. A uniform statement over
program codes uses a bounded evaluator $\operatorname{run}(e,x,t)$, defined
by primitive recursion on $t$, together with a supplied PR clock and its
termination proof. There is no total PR universal evaluator for all PR
programs: if $U(e,x)$ were one, $U(n,n)+1$ would be a PR function and its own
index would give a contradiction. This elementary diagonal argument is a
useful guard against an otherwise tempting but invalid encoding shortcut.

\begin{definition}[A quantitative presentation]\label{def:presentation}
A presentation supplies a collection of finite codes, a PR admissibility
test, a PR common-refinement operation where needed, and an exactly
computable rational pseudodistance $d$ on finite codes after refinement.
The pseudometric laws and refinement invariance are proved from the
finite definitions. A regular name is a fixed PR sampler $a(n)$ satisfying
\begin{equation}\label{eq:regular-name}
 d(a(n),a(m))\le2^{-n}+2^{-m}\qquad(n,m\in\N).
\end{equation}
Two names are equivalent when their distance tends to zero with the given
regular estimates; equivalently, after a harmless precision shift, their
cross-distances are bounded by $2^{-n}+2^{-m}$ for all $n,m$.
\end{definition}

The usual metric completion is one semantic interpretation of these names.
The arithmetic implementation keeps the samplers and the estimates and
need not construct a quotient set. A fixed sampler's regularity is a
first-order formula universally quantifying $n,m$, with a quantifier-free
matrix. Subsequent arguments may use that formula as an assumption; their
induction invariants remain quantifier-free.

If finite distances themselves have certified rational upper and lower
approximations, the same construction works with extra guard bits. The
applications here deliberately use exact rational norms on simple real
functions and finite matrices whenever possible. Complex quantities can be
represented by two real codes and bounded by the rational norm
$|a+ib|_\square=|a|+|b|$; Euclidean-radius guards can instead be checked by
squaring nonnegative rational bounds. These equivalent norm presentations
are connected by explicit inequalities, not identified literally.

\begin{lemma}[Diagonal completion with a supplied clock]\label{lem:diagonal}
Suppose $x_j$ are presented points with jointly PR names, and $N$ is a
nondecreasing PR function such that
\[
 j,\ell\ge N(p)\quad\Longrightarrow\quad d(x_j,x_\ell)\le2^{-p}.
\]
Let $b(p)$ be a finite approximation to $x_{N(p+2)}$ within $2^{-p-2}$.
Then $b$ is a regular name. Its limit $x$ satisfies
$d(x_j,x)\le2^{-p}$ for $j\ge N(p)$.
\end{lemma}
\begin{proof}
For $p\le q$, the common Cauchy estimate gives
\[
 d(b(p),b(q))\le2^{-p-2}+2^{-p-2}+2^{-q-2}
 \le2^{-p}+2^{-q}.
\]
For fixed $j\ge N(p)$, compare $x_j$ with $b(q)$ for $q$ large enough that
$N(q+2)\ge N(p)$. The bound is $2^{-p}+2^{-q-2}$; the named limit removes
the final term. This last sentence is arithmetic shorthand for all its
rational-precision instances. To prove any strict enlargement of $2^{-p}$,
choose a larger explicit $q$; no choice of an unbounded limit index occurs.
Replacing $p$ by $p+1$ supplies a direct weak $2^{-p}$ output bound if a
one-sided name convention is desired. All samplers are compositions of
the supplied PR programs.
\end{proof}

\subsection{Norms are part of the interface}

For a finite probability vector $w=(w_i)_{i<n}$, $w_i\ge0$,
$\sum_iw_i=1$, define
\[
 I_w(f)=\sum_{i<n}w_i f_i,\qquad
 \norm f_{1,w}=\sum_iw_i|f_i|,\qquad
 \norm f_\infty=\max_i|f_i|.
\]
The maximum over an empty list is defined to be zero, although a normalized
probability vector cannot have empty support. Refinement repeats values
on child cells and splits mass, preserving the first two expressions.
For a real matrix use the row-sum norm
$\norm A_\infty=\max_i\sum_j|A_{ij}|$, including dimension-zero cases with
the same convention. All these are PR rational functions.

An operation's contract names both the input and output error. Integration
is $L^1\to|\cdot|$ nonexpansive; taking a supremum is
$\|\cdot\|_\infty\to|\cdot|$ nonexpansive; multiplication is
$L^\infty\times L^1\to L^1$ continuous with a supplied bounded-factor
constant. These distinctions are logically significant.

\begin{proposition}[Safe elementary transfers]\label{prop:safe-transfers}
On a finite probability algebra, and hence on the indicated completions,
\begin{align}
 |I(f)-I(g)|&\le\norm{f-g}_1,\label{eq:int-lipschitz}\\
 \norm{fg-\widetilde f\widetilde g}_1
 &\le B\norm{g-\widetilde g}_1+C\norm{f-\widetilde f}_1,
 \label{eq:product-lipschitz}
\end{align}
provided $|f|\le B$ and $|\widetilde g|\le C$ almost everywhere.
Also $|\sup f-\sup g|\le\norm{f-g}_\infty$ when those suprema are named.
\end{proposition}
\begin{proof}
The first inequality is the triangle inequality for a finite weighted sum.
For the second write $fg-\widetilde f\widetilde g
=f(g-\widetilde g)+(f-\widetilde f)\widetilde g$ and apply the pointwise
triangle inequality before summation. For a supremum, $f(x)\le
g(x)+\norm{f-g}_\infty$ gives one direction, and interchanging $f,g$ gives
the other. Completion preserves inequalities between uniformly continuous
operations by comparing finite approximants with arbitrarily small explicit
errors. The bounded-factor hypotheses in the product estimate must be
preserved or re-certified; $L^1$ names alone do not supply them.
\end{proof}

There is no general transfer from $L^1$ error to uniform error. On $[0,1]$,
$f_n=\mathbf1_{[0,2^{-n}]}$ has $\norm{f_n}_1=2^{-n}$ and
$\norm{f_n}_\infty=1$. Even continuous tent functions give the same
obstruction. Likewise, finite-support atomic measures do not approximate
Lebesgue measure in total variation: their finite support has mass one
for the atomic measure and zero for Lebesgue measure. Finite cell
\emph{observations} of a measure and atomic approximations for Lipschitz
test functions are therefore separate representations. The later category
and model-theory sections use the representation appropriate to each claim.

\subsection{Weighted error and efficient demand propagation}

\begin{lemma}[Mass and value perturbation]\label{lem:weighted}
Let $w_i\ge0$, let $v_i$ be rational, and suppose
$|a_i-b_i|\le e_i$ and $|b_i|\le B$. Then
\begin{equation}\label{eq:weighted}
 \left|\sum_i w_i a_i-\sum_i v_i b_i\right|
 \le\sum_iw_i e_i+B\sum_i|w_i-v_i|.
\end{equation}
If $w=v$ is a probability vector and all $e_i\le e$, the bound is $e$.
\end{lemma}
\begin{proof}
Subtract and add $\sum_iw_ib_i$. The difference is
$\sum_iw_i(a_i-b_i)+\sum_i(w_i-v_i)b_i$.
Apply the triangle inequality to each summand, then use the two supplied
pointwise bounds. This proof works for zero weights and does not divide by
any of them. It also makes explicit which vector weights each local error.
\end{proof}

The repository already has an internal rational error-DAG producer and a
uniform precision schedule \cite{repo-dag,repo-dag-schedule}. A useful
extension is to retain the actual probability coefficients and the norm
type of every edge. The following standard finite calculation explains why.

\begin{proposition}[Backward influence coefficients]\label{prop:adjoint}
For a topologically ordered finite computation suppose its errors satisfy
\[
 E_v\le\delta_v+\sum_{u<v}a_{vu}E_u,\qquad a_{vu}\ge0.
\]
For a selected output $o$, initialize $\lambda_o=1$ and other coefficients
to zero. Process vertices in decreasing order, adding
$\lambda_v a_{vu}$ to $\lambda_u$ along every incoming edge. Then
\[
 E_o\le\sum_v\lambda_v\delta_v.
\]
Consequently any nonnegative local budgets with
$\sum_v\lambda_v\delta_v\le2^{-p}$ certify the output precision.
\end{proposition}
\begin{proof}
Begin with $E_o$. Substitution of the inequality for a vertex replaces its
coefficient times $E_v$ by that coefficient times $\delta_v$ plus the
weighted predecessor errors. Nonnegativity preserves the direction of the
inequality. After each vertex has been processed once, no $E_v$ remains.
The accumulated coefficient of each $\delta_v$ is exactly the value
computed by the reverse pass. This proves the displayed bound, including
graphs with shared subexpressions. A bounded reverse loop and rational
addition/multiplication implement the pass by primitive recursion.
\end{proof}

For integration, $a_{vu}$ are the actual masses. A term of mass $2^{-20}$
can tolerate much larger absolute error than one of mass $1/2$.
Zero-influence terms need no numerical refinement once their range and
well-formedness are established. This is a mathematical error reduction,
not a claim that every resulting algorithm is polynomial-time. Tail
moduli, matrix dimensions, and compact-net sizes may still be enormous.
Max-type rules can be kept as maxima or safely overestimated by sums;
using a sum changes the cost bound but not correctness.

\subsection{Why the finite arguments fit quantifier-free induction}

For \eqref{eq:weighted}, fix a common positive denominator for all finite
input values. After clearing it, let $S_j,T_j,U_j,V_j$ be the relevant
integer prefix sums. A convenient induction invariant is
\[
 -U_j-BV_j\le S_j-T_j\le U_j+BV_j,
\]
together with the stored positivity conditions. This is quantifier-free
in the codes and $j$. At the successor step add the single new weighted
inequality. The zero case is literal equality. The normalization step uses
positive-denominator order preservation already available in the
repository's rational arithmetic. The variable-length sum is thus an
internal numerical function, not an externally expanded expression for
one fixed dimension.

A bounded universal condition is represented by a PR conjunction fold:
$\operatorname{all}_P(0)=1$ and
$\operatorname{all}_P(j+1)=\operatorname{and}
(\operatorname{all}_P(j),P(j))$, where $P$ is a decidable numerical test.
Proofs about a whole input table can use the equality
$\operatorname{all}_P(n)=1$ as a quantifier-free premise and derive its
individual bounded instances. In every main construction below we give
either the corresponding prefix invariant or the finite algebraic identity
that replaces it. There is no rule saying that an arbitrary true numerical
claim is automatically provable in iPRA.
